\section{Model-assisted follow-up and second experiment}
\label{sec:model-followup}
\textbf{Status: completed CPU preparation and audit; physical follow-up unexecuted.}
This separate study asks whether progress-aware control and qualified short-horizon
prediction can improve task continuation without paying for a broad physical
search. Its CPU implementation does not change the frozen comparison above.

\subsection{Implemented model and control interfaces}
A five-member episode-bootstrap ensemble predicts the next 44 compact features
from the current and previous features, a history-valid flag and eight known
executed eight-dimensional commands. Its 153 inputs feed ridge regressors with
32 fixed nonlinear random features. Separate heads estimate a new violation
within the prefix and five task predicates: closed lid, holding, goal
intersection, uprightness and support. A readout evaluated on actual future
features separates perception/readout error from imagined-dynamics error.
This compact transition prototype does not generate future RGB observations or
simulate a complete closed-loop VLA episode. A prediction about a decoded prefix
is a surrogate, not measured fitness.

A prospective controller augments the 44 causal inputs with five learned beliefs.
Its $5\times49$ weights produce four arm-noise coefficients and one gripper-noise
coefficient, with coefficient persistence. Prototype settings are arm scale .35,
gripper scale 2.0 and persistence .8; they are not physically validated magnitudes.
Native Gaussian noise is retained. Privileged simulator predicates cannot silently
replace the unqualified learned beliefs. A separate bounded joint-target residual
provides a prospective representation control; its source-matched CPU clipping
wrapper still requires actual runtime limits and physical identity validation.

The action-changing filter accepts a supported low-risk prefix, substitutes an
eligible alternative, or abstains. Proposed, filtered and actually executed
commands remain distinct. Candidate allocation scores the surviving prefix and
shares predictions for identical filtered prefixes at the same factual anchor.
Absent qualification, allocation preserves abstention and cannot authorize
execution. Decoder requests for discarded proposals remain charged.

The ROSARL interface implements its practical reward/value-extrema penalty rule,
requiring a declared external value estimator and one penalty frozen before
candidate comparison~\cite{rosarl}. It is not a trained task policy or a complete
ROSARL baseline. The shortest-path theory with unsafe absorbing outcomes does
not certify this finite-horizon partially observed system. SafeDreamer already
combines world-model planning and safety objectives~\cite{safedreamer}; this
combination alone is not the proposed contribution.

\subsection{Completed factual audit and failed gate}
The audit protocol was frozen before fitting. Eight shared-initialization
episodes on tasks 0019/0009 supply 1,992 fitting transitions; one development
episode on task 0017 supplies 249 calibration transitions, and one on task 0013
supplies 83 diagnostic transitions. Splits are by whole task, all previously
exposed. Independent checking outcomes are excluded. Final prefixes without a
following encoded observation are omitted, and risk fitting excludes prefixes
starting after a latched violation. These ten existing episodes cost 18,666
original acquisition controls, retained as reused information.

\begin{table}[t]
\caption{Separate CPU model audit on one diagnostic episode. Values are feature
RMSE; rows are correlated measurements within that episode. The predeclared gate
requires a model/persistence ratio at most .90 in every block. All three fail.}
\label{tab:model-audit}
\begin{tabular}{lrrr}
\toprule
Feature block & Model & Persistence & Ratio \\
\midrule
Robot/history & .16805 & .08516 & 1.973 \\
Overview & .04381 & .02848 & 1.538 \\
Wrist & .06380 & .05481 & 1.164 \\
\bottomrule
\end{tabular}
\end{table}

Actual-feature readout MAE is .191 for closure, .209 for holding, .662 for
uprightness and .210 for support, above the .10 gate. The fitting set contains
positive closure and holding in only one episode each, no in-goal positive
episode, and no safe completion. Its observed closure follows a violation.
Calibration has one episode against the minimum 20. The conservative fallback
therefore rejects every prefix; its 100\% negative-prefix alarm rate is not the
false-alarm rate of a qualified classifier. Counterfactual-action coverage
would still require justification after adequate factual calibration.

The audit consumes 27,412 training-row uses and 5,395 model-head row queries,
with zero new physical episodes, simulator controls, VLA calls or encoder images.
Sixteen focused CPU tests validate the initial assembled interfaces; later
allocation/residual and clipping-wrapper suites add 12 and 15 checks, respectively.
These checks establish software behaviour, not manipulation success. The failed
prediction/coverage gate excludes this model from physical action selection.

\subsection{Proposed small continuation screen}
Before further model-guided optimization, test whether a known closure behaviour
can hand control back for body grasping and transport. Freeze one controller
chosen using exposed training data, its prior-information charge, task, fresh
seed, execution order and every switching parameter before the first outcome.
Use four references on one development case: native only, closure only, a fixed
clock handoff, and a privileged supported-lid-closure handoff. The privileged
reference is an engineering feasibility control, with additional actor
information; it cannot establish equal-information learned-method superiority.

The proposed trigger requires three consecutive post-control observations with
the jar closed and supported and no preceding violation. Initial observation
zero does not count. Losing closure/support before the next eight-control query
cancels qualification. Switching occurs at that query and remains latched to
native execution, without resetting its Gaussian RNG. A new trial resets the
phase history; partial failures are charged and cannot be resumed as valid full
episodes. Source pins and per-query branch/RNG receipts need a new validity
reader; the original constant-genotype verifier is insufficient. This contract
has CPU checks, not physical RNG-continuity or continuation evidence.

The first screen is capped at four full starts, including partial and failed
starts. If all are valid, it evaluates the four references once each. Expand
by a separately frozen stage on a second case only if a
prespecified safe ordered body-grasp or completion criterion passes. A final
native-versus-selected-nonprivileged comparison may use four additional full
episodes on fresh seeds: at most 12 valid episodes and two charged infrastructure
replacements, 14 starts total. A valid timeout or adverse outcome is not replaced.
Four full-horizon starts require at most 8,000 policy controls and 1,000 native
policy queries, including partial/invalid work within that screen. Initialization,
settling, encoders, partial failures,
allocated time and preservation require their own numerical ceilings in the
execution-owned protocol. No extra decoder sweep is needed for this handoff
screen. The original study's protected sixteen checking episodes remain separate.

This is a prospective reduction of the earlier 16-valid-plus-two-replacement
capability proposal, not an additional allowance. Execution must freeze the
complete staged protocol and verify resource/deadline capacity before launch.
Even a successful small screen cannot supply the minimum 20 independent
calibration episodes for the learned filter. Failure calls for diagnosis rather
than a larger optimizer sweep and does not establish irreversibility.

\subsection{Scientific comparison after capability and model repair}
A later funded comparison must separate model-free and model-assisted evolution,
matched random allocation/search, ordinary task-policy learning with the same
filter, and native, filtered-native and strong fixed references. Give competitors
the same observations, training labels, prior data, model access and total
interaction budget. Progress-aware training must retain requested safe completion
as the outcome and remove progress fitness after a violation. Filter, ROSARL and
evolutionary retention each need their own ablation. Replicated optimizers,
independent confirmation and frozen-controller family transfer remain necessary
before asserting the intended GECCO contribution.
